\documentclass[usenames,dvipsnames]{beamer}
%\documentclass[handout]{beamer}
%\usepackage{etex}
\usepackage[scaled=0.92]{helvet}	% set Helvetica as the sans-serif font
\renewcommand{\rmdefault}{ptm}		% set Times as the default text font
\usepackage{xcolor}
%\usepackage{definice}
\usepackage{mathptmx}
\usepackage{amsmath,amssymb}
\usepackage{mathtools}
\usepackage{tikz}
% Bold latin letters
\newcommand{\bfa}{\ensuremath{\mathbf{a}}}
\newcommand{\bfb}{\ensuremath{\mathbf{b}}}
\newcommand{\bfc}{\ensuremath{\mathbf{c}}}
\newcommand{\bfd}{\ensuremath{\mathbf{d}}}
\newcommand{\bfe}{\ensuremath{\mathbf{e}}}
\newcommand{\bff}{\ensuremath{\mathbf{f}}}
\newcommand{\bfg}{\ensuremath{\mathbf{g}}}
\newcommand{\bfh}{\ensuremath{\mathbf{h}}}
\newcommand{\bfi}{\ensuremath{\mathbf{i}}}
\newcommand{\bfj}{\ensuremath{\mathbf{j}}}
\newcommand{\bfk}{\ensuremath{\mathbf{k}}}
\newcommand{\bfl}{\ensuremath{\mathbf{l}}}
\newcommand{\bfm}{\ensuremath{\mathbf{m}}}
\newcommand{\bfn}{\ensuremath{\mathbf{n}}}
\newcommand{\bfo}{\ensuremath{\mathbf{o}}}
\newcommand{\bfp}{\ensuremath{\mathbf{p}}}
\newcommand{\bfq}{\ensuremath{\mathbf{q}}}
\newcommand{\bfr}{\ensuremath{\mathbf{r}}}
\newcommand{\bfs}{\ensuremath{\mathbf{s}}}
\newcommand{\bft}{\ensuremath{\mathbf{t}}}
\newcommand{\bfu}{\ensuremath{\mathbf{u}}}
\newcommand{\bfv}{\ensuremath{\mathbf{v}}}
\newcommand{\bfw}{\ensuremath{\mathbf{w}}}
\newcommand{\bfx}{\ensuremath{\mathbf{x}}}
\newcommand{\bfy}{\ensuremath{\mathbf{y}}}
\newcommand{\bfz}{\ensuremath{\mathbf{z}}}

\newcommand{\bfA}{\ensuremath{\mathbf{A}}}
\newcommand{\bfB}{\ensuremath{\mathbf{B}}}
\newcommand{\bfBu}{\ensuremath{\mathbf{B}_{\mathrm{u}}}}
\newcommand{\bfBd}{\ensuremath{\mathbf{B}_{\mathrm{d}}}}
\newcommand{\bfBf}{\ensuremath{\mathbf{B}_{\mathrm{f}}}}
\newcommand{\bfBw}{\ensuremath{\mathbf{B}_{\mathrm{w}}}}
\newcommand{\bfC}{\ensuremath{\mathbf{C}}}
\newcommand{\bfD}{\ensuremath{\mathbf{D}}}
\newcommand{\hbfD}{\ensuremath{\hat{\mathbf{D}}}}
\newcommand{\bfDu}{\ensuremath{\mathbf{D}_{\mathrm{u}}}}
\newcommand{\bfDv}{\ensuremath{\mathbf{D}_{\mathrm{v}}}}
\newcommand{\bfDd}{\ensuremath{\mathbf{D}_{\mathrm{d}}}}
\newcommand{\bfDf}{\ensuremath{\mathbf{D}_{\mathrm{f}}}}
\newcommand{\bfE}{\ensuremath{\mathbf{E}}}
\newcommand{\bfF}{\ensuremath{\mathbf{F}}}
\newcommand{\bfG}{\ensuremath{\mathbf{G}}}
\newcommand{\bfH}{\ensuremath{\mathbf{H}}}
\newcommand{\bfI}{\ensuremath{\mathbf{I}}}
\newcommand{\bfJ}{\ensuremath{\mathbf{J}}}
\newcommand{\bfK}{\ensuremath{\mathbf{K}}}
\newcommand{\bfL}{\ensuremath{\mathbf{L}}}
\newcommand{\bfM}{\ensuremath{\mathbf{M}}}
\newcommand{\bfN}{\ensuremath{\mathbf{N}}}
\newcommand{\bfO}{\ensuremath{\mathbf{O}}}
\newcommand{\bfP}{\ensuremath{\mathbf{P}}}
\newcommand{\bfPf}{\ensuremath{\mathbf{P}}}
\newcommand{\bfPp}{\ensuremath{\mathbf{P}'}}
\newcommand{\bfQ}{\ensuremath{\mathbf{Q}}}
\newcommand{\bfR}{\ensuremath{\mathbf{R}}}
\newcommand{\bfS}{\ensuremath{\mathbf{S}}}
\newcommand{\bfT}{\ensuremath{\mathbf{T}}}
\newcommand{\bfU}{\ensuremath{\mathbf{U}}}
\newcommand{\bfV}{\ensuremath{\mathbf{V}}}
\newcommand{\bfW}{\ensuremath{\mathbf{W}}}
\newcommand{\bfX}{\ensuremath{\mathbf{X}}}
\newcommand{\bfY}{\ensuremath{\mathbf{Y}}}
\newcommand{\bfZ}{\ensuremath{\mathbf{Z}}}

% Bold latin letters with bar
\newcommand{\bbfx}{\ensuremath{\bar{\mathbf{x}}}}
\newcommand{\bbfy}{\ensuremath{\bar{\mathbf{y}}}}
\newcommand{\bbfz}{\ensuremath{\bar{\mathbf{z}}}}

\newcommand{\bbfA}{\ensuremath{\bar{\mathbf{A}}}}
\newcommand{\bbfG}{\ensuremath{\bar{\mathbf{G}}}}
\newcommand{\bbfI}{\ensuremath{\bar{\mathbf{I}}}}
\newcommand{\bbfZ}{\ensuremath{\bar{\mathbf{Z}}}}

% Bold latin letters with hat
\newcommand{\hbfd}{\ensuremath{\hat{\mathbf{d}}}}
\newcommand{\hbfv}{\ensuremath{\hat{\mathbf{v}}}}
\newcommand{\hbfx}{\ensuremath{\hat{\mathbf{x}}}}
\newcommand{\hbfy}{\ensuremath{\hat{\mathbf{y}}}}
\newcommand{\hbfz}{\ensuremath{\hat{\mathbf{z}}}}


% Bold latin letters with dot
\newcommand{\dbfx}{\ensuremath{\dot{\mathbf{x}}}}

% Bold latin letters with tilde
\newcommand{\tbfx}{\ensuremath{\tilde{\mathbf{x}}}}


\newcommand{\bffu}[1][,k]{\ensuremath{\mathbf{f}_{\mathrm{u}#1}}}
\newcommand{\bffy}[1][,k]{\ensuremath{\mathbf{f}_{\mathrm{y}#1}}}


\newcommand{\Lk}{\ensuremath{L_{k}\left(\bfx_{k},\bfmu_{k},\bfu_{k},\bfd_{k}\right)}}
\newcommand{\Lok}{\ensuremath{L_{k}\left(\bfx_{0}^{k},\bfmu_{0}^{k},\bfu_{0}^{k},\bfd_{0}^{k}\right)}}
\newcommand{\Lkd}{\ensuremath{L_{k}^{\text{d}}\left(\bfmu_{k},\bfd_{k}\right)}}
\newcommand{\Lkds}[1][k]{\ensuremath{L_{#1}^{\text{d}}\left(\mu_{#1},d_{#1}\right)}}
\newcommand{\Lkc}{\ensuremath{L_{k}^{\text{c}}\left(\bfx_{k},\bfu_{k}\right)}}
\newcommand{\Lkcs}{\ensuremath{L_{k}^{\text{c}}\left(x_{k},u_{k}\right)}}
\newcommand{\Lmaxd}{\ensuremath{L_{\mathrm{max}}^{\text{d}}}}
\newcommand{\Lmaxc}{\ensuremath{L_{\mathrm{max}}^{\text{c}}}}
\newcommand{\Ld}{\ensuremath{L^{\text{d}}\left(\mu,d\right)}}
\newcommand{\Lkdparam}[1][d_{k}]{\ensuremath{L^{\mathrm{d}}(\mu_{k},#1)}}

\newcommand{\bfIdent}{\ensuremath{\mathbf{I}}}
\newcommand{\bfZero}{\ensuremath{\mathbf{0}}}
\newcommand{\bfOne}{\ensuremath{\mathbf{1}}}
\newcommand{\hmu}{\ensuremath{\hat{\mu}}}



\newcommand{\sub}[2][S]{\ensuremath{\mathrm{#1}{#2}}}

% Set of natural numbers
\newcommand{\bbN}{\ensuremath{\mathbb{N}}}

% Set of integer numbers
\newcommand{\bbZ}{\ensuremath{\mathbb{Z}}}

% Set of real numbers
\newcommand{\bbR}{\ensuremath{\mathbb{R}}}

% Set of complex numbers
\newcommand{\bbC}{\ensuremath{\mathbb{C}}}

% Set of discrete time steps
\newcommand{\bbT}{\ensuremath{\mathbb{T}}}


% Calligraphic letters
\newcommand{\mcA}{\ensuremath{\mathcal{A}}}
\newcommand{\mcB}{\ensuremath{\mathcal{B}}}
\newcommand{\mcC}{\ensuremath{\mathcal{C}}}
\newcommand{\mcD}{\ensuremath{\mathcal{D}}}
\newcommand{\mcE}{\ensuremath{\mathcal{E}}}
\newcommand{\mcF}{\ensuremath{\mathcal{F}}}
\newcommand{\mcG}{\ensuremath{\mathcal{G}}}
\newcommand{\mcH}{\ensuremath{\mathcal{H}}}
\newcommand{\mcI}{\ensuremath{\mathcal{I}}}
\newcommand{\mcJ}{\ensuremath{\mathcal{J}}}
\newcommand{\mcK}{\ensuremath{\mathcal{K}}}
\newcommand{\mcL}{\ensuremath{\mathcal{L}}}
\newcommand{\mcM}{\ensuremath{\mathcal{M}}}
\newcommand{\mcN}{\ensuremath{\mathcal{N}}}
\newcommand{\mcO}{\ensuremath{\mathcal{O}}}
\newcommand{\mcP}{\ensuremath{\mathcal{P}}}
\newcommand{\mcQ}{\ensuremath{\mathcal{Q}}}
\newcommand{\mcR}{\ensuremath{\mathcal{R}}}
\newcommand{\mcS}{\ensuremath{\mathcal{S}}}
\newcommand{\mcT}{\ensuremath{\mathcal{T}}}
\newcommand{\mcU}{\ensuremath{\mathcal{U}}}
\newcommand{\mcV}{\ensuremath{\mathcal{V}}}
\newcommand{\mcW}{\ensuremath{\mathcal{W}}}
\newcommand{\mcX}{\ensuremath{\mathcal{X}}}
\newcommand{\mcY}{\ensuremath{\mathcal{Y}}}
\newcommand{\mcZ}{\ensuremath{\mathcal{Z}}}

% Math script letters
\newcommand{\mkZ}{\ensuremath{\mathscr{Z}}}
\newcommand{\mkF}{\ensuremath{\mathscr{F}}}

% Definition of math operators
\DeclareMathOperator*{\mean}{E}				% Expectation operator
\DeclareMathOperator*{\cov}{cov}			% Covariance operator
\DeclareMathOperator*{\var}{var}			% Variance operator
\DeclareMathOperator*{\sign}{sign}			% Signum operator
\DeclareMathOperator*{\tr}{tr}				% Trace operator
\DeclareMathOperator*{\rank}{rank}			% Rank operator
\DeclareMathOperator*{\diff}{d}				% Differential operator
\DeclareMathOperator*{\argmin}{arg\,min}	% Argmin operator
\DeclareMathOperator*{\argmax}{arg\,max}	% Argmax operator
\DeclareMathOperator*{\range}{range}		% Range operator
\DeclareMathOperator*{\nullsp}{null}		% Range operator
\DeclareMathOperator*{\vspan}{span}			% Span operator
\DeclareMathOperator*{\dom}{dom}			% Domain oprator operator
\DeclareMathOperator*{\vect}{vec}			% Vectorization operator
\DeclareMathOperator*{\prst}{P}				% Probability operator

\newcommand{\indfun}{\ensuremath{\mathds{1}}} % Indicator function requires \usepackage{dsfont}
\newcommand{\iu}{\ensuremath{i}}        % Imaginary unit
\providecommand{\e}{\ensuremath{e}}        % Euler's number
\providecommand{\norm}[1]{\lVert#1\rVert}
\newcommand{\T}{\ensuremath{\text{T}}} % Transpose
\newcommand{\Ts}{\ensuremath{T_{\mathrm{s}}}}
\newcommand{\qo}{q^{-1}}
\newcommand{\SISO}{\textsf{SISO}}
\newcommand{\MIMO}{\textsf{MIMO}}
\newcommand{\kf}{\ensuremath{k_{\mathrm{f}}}}
\newcommand{\ka}{\ensuremath{k_{\mathrm{a}}}}
\newcommand{\kfh}{\ensuremath{\hat{k}_{\mathrm{f}}}}

% Bold Greek letters
\newcommand{\bfalpha}{\ensuremath{\boldsymbol{\alpha}}}
\newcommand{\bfbeta}{\ensuremath{\boldsymbol{\beta}}}
\newcommand{\bfdelta}{\ensuremath{\boldsymbol{\delta}}}
\newcommand{\bfDelta}{\ensuremath{\boldsymbol{\Delta}}}
\newcommand{\bfmu}{\ensuremath{\boldsymbol{\mu}}}
\newcommand{\bflambda}{\ensuremath{\boldsymbol{\lambda}}}
\newcommand{\bfLambda}{\ensuremath{\boldsymbol{\Lambda}}}
\newcommand{\bfomega}{\ensuremath{\boldsymbol{\omega}}}
\newcommand{\bftau}{\ensuremath{\boldsymbol{\tau}}}
\newcommand{\bftheta}{\ensuremath{\boldsymbol{\theta}}}
\newcommand{\bfrho}{\ensuremath{\boldsymbol{\rho}}}
\newcommand{\bfeta}{\ensuremath{\boldsymbol{\eta}}}
\newcommand{\bfpsi}{\ensuremath{\boldsymbol{\psi}}}
\newcommand{\bfpi}{\ensuremath{\boldsymbol{\pi}}}
\newcommand{\bfnu}{\ensuremath{\boldsymbol{\nu}}}
\newcommand{\bfsigma}{\ensuremath{\boldsymbol{\sigma}}}
\newcommand{\bfgamma}{\ensuremath{\boldsymbol{\gamma}}}
\newcommand{\bfchi}{\ensuremath{\boldsymbol{\chi}}}
\newcommand{\bfOmega}{\ensuremath{\boldsymbol{\Omega}}}
\newcommand{\bfGamma}{\ensuremath{\boldsymbol{\Gamma}}}
\newcommand{\bfSigma}{\ensuremath{\boldsymbol{\Sigma}}}
\newcommand{\bfSigmav}{\ensuremath{\boldsymbol{\Sigma}_{\mathrm{v}}}}
\newcommand{\bfSigmaw}{\ensuremath{\boldsymbol{\Sigma}_{\mathrm{w}}}}
\newcommand{\bfvarphi}{\ensuremath{\boldsymbol{\varphi}}}
\newcommand{\bfphi}{\ensuremath{\boldsymbol{\phi}}}
\newcommand{\bfPhi}{\ensuremath{\boldsymbol{\Phi}}}
\newcommand{\bfepsilon}{\ensuremath{\boldsymbol{\epsilon}}}
\newcommand{\bfxi}{\ensuremath{\boldsymbol{\xi}}}
\newcommand{\bfzeta}{\ensuremath{\boldsymbol{\zeta}}}

% Bold Greek letters with hat
\newcommand{\hbfxi}{\ensuremath{\hat{\boldsymbol{\xi}}}}

% Special symbols
\newcommand{\Ak}{\ensuremath{\bfA_{\mu_{k}}}}
\newcommand{\Bk}{\ensuremath{\bfB_{\mu_{k}}}}
\newcommand{\Ck}{\ensuremath{\bfC_{\mu_{k}}}}
\newcommand{\Gk}{\ensuremath{\bfG_{\mu_{k}}}}
\newcommand{\Hk}{\ensuremath{\bfH_{\mu_{k}}}}

% Software
\RequirePackage{xspace}
\newcommand{\matlab}{\mbox{\textsc{Matlab}}\xspace}
\newcommand{\yalmip}{\mbox{\textsc{Yalmip}}\xspace}
\newcommand{\cplex}{\mbox{\textsc{CPLEX}}\xspace}
\newcommand{\snopt}{\mbox{\textsc{SNOPT}}\xspace}
\newcommand{\bmibnb}{\mbox{\texttt{bmibnb}}\xspace} 
\usepackage{slashbox}
\usepackage{booktabs,hhline}
%\usepackage{xcolor}%,colortbl}
\usepackage{multirow}
\usepackage{adjustbox}
\usepackage{siunitx}
%\usetikzlibrary{arrows,calc,fadings,positioning,shapes.arrows,decorations.pathreplacing,calc}
\usetikzlibrary{decorations.markings,shadows.blur,fit,arrows,calc,fadings,positioning,shapes.arrows,shadows,backgrounds,fit}
\usepackage[beamer,customcolors]{hf-tikz}
\hfsetfillcolor{blue!20}
\hfsetbordercolor{blue!50}
%\usepackage{pgf}
%\usepackage{pgfplots}
%\usetikzlibrary{shapes.arrows,decorations.pathmorphing,positioning,datavisualization.formats.functions,calc,fadings,decorations.pathreplacing,mindmap,trees,positioning,arrows,shapes,shadows,arrows,backgrounds,fit,plotmarks}
\let\Bbbk\relax
\let\jmath\relax
\let\hbar\relax
\usepackage[subscriptcorrection,slantedGreek,nofontinfo,mtphrb,mtpscr]{mtpro2_fix}
%\usepackage{lmodern}
\usepackage{beamerhighlight}

\setbeamertemplate{navigation symbols}{}
\mode<beamer>{
  \setbeamertemplate{background canvas}[vertical shading][bottom=black!15,middle=black!5,top=black!15]
  \usetheme{metropolis}
  \definecolor{UBCblue}{rgb}{0.04706, 0.13725, 0.46667} % UBC Blue (primary)
  \usecolortheme[named=UBCblue]{structure}
}
\mode<handout>{
  \usetheme{default}
}
\setbeamercovered{transparent}

\usepackage{algorithm2e}
%----------------------------------
% set algorithm2e styles
%----------------------------------
% change algorithm font size
\SetAlFnt{\footnotesize}

% change algorithm caption style
\newcommand{\xAlCapSty}[1]{\small\sffamily\bfseries\MakeUppercase{#1}}
\SetAlCapSty{xAlCapSty}

% comment style (algorithms)
\newcommand{\xCommentSty}[1]{\scriptsize\ttfamily\textcolor{blue}{#1}}
\SetCommentSty{xCommentSty}

% change line number style
\newcommand\mynlfont[1]{\scriptsize\sffamily{#1}}
\SetNlSty{mynlfont}{}{} 

% don't print semicolon
\DontPrintSemicolon

% ruled algorithm
\RestyleAlgo{algoruled}

\makeatletter
\newcommand{\removelatexerror}{\let\@latex@error\@gobble}
\makeatother
\tikzset{
  every overlay node/.style={
    %draw=black,fill=white,rounded corners,
    anchor=north west, inner sep=0pt,
  },
}

\graphicspath{{img/}}

\setbeamertemplate{footline}
{%
  \hfill%
  \includegraphics[width=1.5cm]{img/zcu_en_cmyk.pdf}
}

%\DeclareMathAlphabet{\pazocal}{OMS}{zplm}{m}{n}
\tikzset{
	sys/.style={rectangle,minimum height=1.1cm,minimum width=2.1cm,rounded corners,draw=black,very thick,align=center},
	curvedlink/.style={to path={let \p1=(\tikztostart.east), \p2=(\tikztotarget.west),\n1= {abs(\y2-\y1)/4} in (\p1) arc(90:-90:\n1) -- ([yshift=2*\n1]\p2) arc (90:270:\n1)}}
}
\tikzstyle{wrlss}=[rectangle,thick,minimum size=1.8cm,draw=gray!80,fill=gray!10,blur shadow     = {shadow blur steps = 15}]
\tikzstyle{mlss}=[rectangle,thick,minimum size=1.8cm,draw=black,fill=white=!90,fill opacity=0.5]
\tikzstyle{wr}=[rectangle,thick,minimum size=1cm,draw=gray!80,fill=gray!20]
\tikzstyle{rr}=[rectangle,thick,minimum size=1cm,draw=gray!80,fill=red!20,blur shadow={shadow blur steps = 15}]
\tikzstyle{yr}=[rectangle,thick,minimum size=1cm,draw=gray!80,fill=yellow!40,blur shadow={shadow blur steps = 15}]
\tikzstyle{pr}=[rectangle,thick,minimum size=1cm,draw=gray!80,fill=purple!60,blur shadow={shadow blur steps = 15}]
\tikzstyle{gr}=[rectangle,thick,minimum size=1cm,draw=gray!80,fill=green!20,blur shadow={shadow blur steps = 15}]
\tikzstyle{grVecArrow} = [decoration={markings,mark=at position 0.2 with {\arrow{triangle 60[reversed,fill=green!20]}}},double distance=1.4pt, shorten <= 5.5pt, preaction = {decorate}, postaction = {draw,line width=1.4pt, green!20,shorten <= 4.5pt},decoration={markings,mark=at position 1 with {\arrow{triangle 60[fill=green!20]}}}, double distance=1.4pt, shorten >= 4.5pt,  preaction = {decorate},  postaction = {draw,line width=1.4pt, green!20,shorten >= 5.3pt}]
\tikzstyle{grVecArrow1} = [decoration={markings,mark=at position 0.07 with {\arrow{triangle 60[reversed,fill=green!20]}}},double distance=1.4pt, shorten <= 5.5pt, preaction = {decorate}, postaction = {draw,line width=1.4pt, green!20,shorten <= 4.5pt},decoration={markings,mark=at position 1 with {\arrow{triangle 60[fill=green!20]}}}, double distance=1.4pt, shorten >= 4.5pt,  preaction = {decorate},  postaction = {draw,line width=1.4pt, green!20,shorten >= 5.3pt}]
\tikzstyle{purVecArrow} = [decoration={markings,mark=at position 0.15 with {\arrow{triangle 60[reversed,fill=purple!20]}}},double distance=1.4pt, shorten <= 5.5pt, preaction = {decorate}, postaction = {draw,line width=1.4pt, purple!20,shorten <= 4.5pt},decoration={markings,mark=at position 1 with {\arrow{triangle 60[fill=purple!20]}}}, double distance=1.4pt, shorten >= 4.5pt,  preaction = {decorate},  postaction = {draw,line width=1.4pt, purple!20,shorten >= 5.3pt}]
\tikzstyle{purVecArrow1} = [decoration={markings,mark=at position 0.1 with {\arrow{triangle 60[reversed,fill=purple!20]}}},double distance=1.4pt, shorten <= 5.5pt, preaction = {decorate}, postaction = {draw,line width=1.4pt, purple!20,shorten <= 4.5pt},decoration={markings,mark=at position 1 with {\arrow{triangle 60[fill=purple!20]}}}, double distance=1.4pt, shorten >= 4.5pt,  preaction = {decorate},  postaction = {draw,line width=1.4pt, purple!20,shorten >= 5.3pt}]
\tikzstyle{whiVecArrow} = [thick,decoration={markings,mark=at position 0.2 with {\arrow[semithick]{triangle 60[reversed,fill=white]}}},double distance=1.4pt, shorten <= 3.4pt, preaction = {decorate}, postaction = {draw,line width=1.4pt, white,shorten <= 4.5pt}, decoration={markings,mark=at position 1 with {\arrow[semithick]{triangle 60[fill=white]}}},  double distance=1.4pt, shorten >= 4.5pt,  preaction = {decorate},  postaction  = {draw,line width=1.4pt, white,shorten >= 5.0pt}]
  
\pgfdeclarelayer{back}
\pgfsetlayers{back,main}

\def\zeroth{\prescript{0\!}{}}
\def\ist{\prescript{1\!}{}}
\def\iind{\prescript{2\!}{}}
\def\iiird{\prescript{3\!}{}}
\def\nth{\prescript{n\!}{}}
\def\ith{\prescript{i\!}{}}
\def\jth{\prescript{j\!}{}}
\def\bnth{\prescript{\bar{n}\!}{}}
\def\Nth{\prescript{N\!}{}}
\def\bffX{\mathbb{X}}
%\def\bffX{\mathds{X}}
%%%%%%%%%%%%%%%%%%%%%%%%%%%%%%%%%%%%%%%%%%%%%%%%%%%%%%%%%%%%%%%%%%%%%%%%%%%%%%%%%%%%%%%%%%%%%%%%%%
%%%%%%%%%%%%%%%%%%%%%%%%%%%%%%%%%%%%%%%%%%%%%%%%%%%%%%%%%%%%%%%%%%%%%%%%%%%%%%%%%%%%%%%%%%%%%%%%%%
\title{Hierarchical Active Fault Diagnosis for Stochastic Large Scale Systems with Coupled Faults}
\author{\textbf{Ond\v{r}ej Straka} and Ivo Pun\v{c}och\'{a}\v{r}\\\small 
\textbf{straka30@kky.zcu.cz}, ivop@kky.zcu.cz}
\institute{University of West Bohemia\\
Pilsen, Czech Republic}
\date{FUSION 2020,\\ July 6--9, 2020}
\begin{document}
\frame{\titlepage}
\begin{frame}\frametitle{Active vs. Passive Fault Diagnosis}
 \begin{tikzpicture}
    %\node [sys,fill=red!15,draw=red!50] (systemp) {System};
    %\node [sys,fill=brown!15,draw=brown!50,minimum width=4cm] at ($(systemp)+(4.5,0)$) (pfd) {Passive Fault Detector};
    %\draw[thick,->] ($(systemp.east)+(0,0)$) -- ($(pfd.west)+(0,0)$) node [above,midway] {Data};
    %\draw[thick,->] ($(pfd.east)+(0,0)$) -- ($(pfd.east)+(1.8,0)$) node [above,midway] {Decision};

    %\begin{pgfonlayer}{back}
    %\draw[dashed,fill=black!5] ($(systemp.north west) + (-3mm,3mm)$) rectangle ($(pfd.south east)+(2.1,-3mm)$);
    %\end{pgfonlayer}
   \begin{scope}
    \node [sys,fill=blue!15,draw=blue!50] (system) {System};
    \node[sys,fill=red!15,draw=red!50,minimum width=4cm] at ($(system)+(4.5,0)$) (afd) {Passive Fault Detector};
    \draw[thick,->] ($(system.east)+(0,0.3)$) -- ($(afd.west)+(0,0.3)$) node [above,midway] {Data};
    \draw[thick,->] ($(afd.east)+(0,0)$) -- ($(afd.east)+(1.8,0)$) node [above,midway] {Decision};
    \begin{pgfonlayer}{back}
      \draw[dashed,fill=black!5] ($(system.north west) + (-3mm,3mm)$) rectangle ($(afd.south east)+(2.1,-3mm)$);
    \end{pgfonlayer}
   \end{scope}
   \begin{scope}[yshift=-4cm]
    \node [sys,fill=blue!15,draw=blue!50] (system) {System};
    \node[sys,fill=red!15,draw=red!50,minimum width=4cm] at ($(system)+(4.5,0)$) (afd) {Active Fault Detector};
    \draw[thick,->] ($(system.east)+(0,0.3)$) -- ($(afd.west)+(0,0.3)$) node [above,midway] {Data};
    \draw[thick,->] ($(afd.east)+(0,0)$) -- ($(afd.east)+(1.8,0)$) node [above,midway] {Decision};
    \draw[thick,->]  ($(afd.west)+(0,-0.3)$)  -- node [below,midway] {Excitation} ($(system.east)+(0,-0.3)$);
    \begin{pgfonlayer}{back}
      \draw[dashed,fill=black!5] ($(system.north west) + (-3mm,3mm)$) rectangle ($(afd.south east)+(2.1,-3mm)$);
    \end{pgfonlayer}
   \end{scope}
  %\only<3>{\end{scope}}

    %\visible<3>{
      %\node[sys,below=2cm of system,fill=red!10,xshift=1cm] (rg) {Residue generator};
      %\node[left=1.6cm of rg] (data) {};
      %\node[sys,right=1.5cm of rg,fill=red!10] (dg) {Decision generator};
      %\node[sys,below=1cm of dg,fill=red!10](ig) {Input signal generator};
      %\node[left=6cm of ig] (excitation) {};
      %\node[right=1.5 cm of dg] (decision) {};
      %\draw[thick,->] (data) -- node[above] {Data} (rg);
      %\draw[thick,->] (rg) -- node[above] {Residue} (dg);
      %\draw[thick,->] (dg) -- node[above] {Decision} (decision);
      %\draw[thick,->] (ig) -- node[above,xshift=-2cm] {Excitation} (excitation);
      %\draw[thick,->] (ig) -| (rg);
      %\draw[thick,->] ($(rg.west) +(-2mm,0)$) -- +(0,-1cm) -| (ig.north);
      %\draw[fill] ($(rg.west) +(-2mm,0)$) circle (1.5pt);
      %\draw[fill] (ig-|rg) circle (1.5pt);
      %\begin{pgfonlayer}{back}
        %\draw<3>[fill=red!15,draw=red!50,thick,rounded corners,very thick] ($(rg.north west) + (-3mm,3mm)$) rectangle ($(ig.south east)+(3mm,-3mm)$);
        %\draw<3>[fill=red!3,draw=red!30] (afd.south west) -- ($(rg.north west) +
          %(-3mm,3mm)$) -- ($(dg.north east) + (5mm,3mm)$) -- (afd.south east);
    %\end{pgfonlayer}}
  \end{tikzpicture} 
\end{frame}
\begin{frame}
  \frametitle{Large Scale System, Multiple Model Framework -- Illustration}
  \vfill
  \begin{center}
    \begin{tikzpicture}[scale=0.7]
      \node (dix1) [wrlss,label={[yshift=-4.4mm,xshift=-.3mm]\tiny subsystem1
        ($\ist{\Sigma}$)}] {};
      \node (dix2) [wrlss,right=1.1cm of
        dix1,label={[yshift=-4.4mm,xshift=-.3mm]\tiny subsystem2
        ($\iind{\Sigma}$)}] {};
      \node (dix3) [wrlss,right=1.1cm of
        dix2,label={[yshift=-4.4mm,xshift=-.3mm]\tiny subsystemN
        ($\Nth{\Sigma}$)}] {};
      \draw[latex-latex,thick] (dix1.35) -- (dix2.145);
      \draw[latex-latex,thick,dotted] (dix2.35) -- (dix3.145);
      \draw[latex-latex,thick] (dix1.north) to [out=10,in=170]  (dix3.north);

      \node (ldi) at ($(dix1.north east)+(-1,.7)$) {LSS ($\Sigma$)};

      \node (m11) [mlss,right=-17mm of
        dix1,yshift=-4mm,label={[yshift=-4.4mm,xshift=-.3mm]\tiny mode1
        ($\ist{\mu}_k=1$)}]  {};
      \node (m12) [mlss,right=-17mm of
        m11,yshift=-4mm,label={[yshift=-4.4mm,xshift=-.3mm]\tiny mode2 ($\ist{\mu}_k=2$)}]  {};
      \node (m13) [mlss,right=-17mm of
        m12,yshift=-4mm,label={[yshift=-4.4mm,xshift=-.3mm]\tiny mode3 ($\ist{\mu}_k=3$)}]  {};

      \node (m21) [mlss,right=-17mm of
        dix2,yshift=-4mm,label={[yshift=-4.4mm,xshift=-.3mm]\tiny mode1
        ($\iind{\mu}_k=1$)}]  {};
      \node (m22) [mlss,right=-17mm of
        m21,yshift=-4mm,label={[yshift=-4.4mm,xshift=-.3mm]\tiny mode2 ($\iind{\mu}_k=2$)}]  {};
      \node (m23) [mlss,right=-17mm of
        m22,yshift=-4mm,label={[yshift=-4.4mm,xshift=-.3mm]\tiny mode3 ($\iind{\mu}_k=3$)}]  {};
      \node (m24) [mlss,right=-17mm of
        m23,yshift=-4mm,label={[yshift=-4.4mm,xshift=-.3mm]\tiny mode4 ($\iind{\mu}_k=4$)}]  {};


      \node (m31) [mlss,right=-17mm of
        dix3,yshift=-4mm,label={[yshift=-4.4mm,xshift=-.3mm]\tiny mode1
        ($\Nth{\mu}_k=1$)}]  {};
      \node (m32) [mlss,right=-17mm of
        m31,yshift=-4mm,label={[yshift=-4.4mm,xshift=-.3mm]\tiny mode2 ($\Nth{\mu}_k=2$)}]  {};
      %\node (m33) [mlss,right=-16mm of
        %m32,yshift=-4mm,label={[yshift=-4.4mm,xshift=-.3mm]\tiny mode3 ($\Nth{\mu}_k=3$)}]  {};
      \begin{pgfonlayer}{back}
        \draw[dashed,blue,fill=blue!5] ($(dix1.north west)+(-.3,1)$) rectangle ($(dix3.south east)+(.5,-.5)$);
      \end{pgfonlayer}
      \end{tikzpicture}
    \end{center}
    \vspace{-3ex}
    \begin{center}
      \scalebox{0.8}{\parbox{\textwidth}{%
          \begin{align*}
            \hspace{-.5em}
            \Sigma:\,
            \underbrace{
              \begin{bmatrix}
                \ist{\bfx}_{k+1}\\
                %\iind{\bfx}_{k+1}\\
                \vdots\\
                \Nth{\bfx}_{k+1}
            \end{bmatrix}}_{\bfx_{k+1}}&\!=\!
            \underbrace{
              \begin{bmatrix}
                \ist{\bff}(\bfx_{k},\ist{\mu}_{k},\ist{\bfu}_{k})\\
                %\iind{\bff}(\bfx_{k},\iind{\mu}_{k},\iind{\bfu}_{k})\\
                \vdots\\
                \Nth{\bff}(\bfx_{k},\Nth{\mu}_{k},\Nth{\bfu}_{k})
            \end{bmatrix}}_{\bff(\bfx_{k},\bfmu_{k},\bfu_{k})}+
            \bfF(\bfmu_k)
            \underbrace{
              \begin{bmatrix}
                \ist{\bfw}_{k}\\
                %\iind{\bfw}_{k}\\
                \vdots\\
                \Nth{\bfw}_{k}
            \end{bmatrix}}_{\bfw_k},\\
            \hspace{-.5em}
            \underbrace{
              \begin{bmatrix}
                \ist{\bfy}_{k}\\
                %\iind{\bfy}_{k}\\
                \vdots\\
                \Nth{\bfy}_{k}
            \end{bmatrix}}_{\bfy_{k}}&=
            \underbrace{
              \begin{bmatrix}
                \ist{\bfh}(\ist{\bfx}_{k},\ist{\mu}_{k})\\
                %\iind{\bfh}(\iind{\bfx}_{k},\iind{\mu}_{k})\\
                \vdots\\
                \Nth{\bfh}(\Nth{\bfx}_{k},\Nth{\mu}_{k})
            \end{bmatrix}}_{\bfh(\bfx_{k},\bfmu_{k})}+
            \bfH(\bfmu_k)
            \underbrace{
              \begin{bmatrix}
                \ist{\bfv}_{k}\\
                %\iind{\bfv}_{k}\\
                \vdots\\
                \Nth{\bfv}_{k}
            \end{bmatrix}}_{\bfv_k},
      \end{align*}}}
    \end{center}
  \end{frame}
\begin{frame}
  \frametitle{Active Fault Diagnosis -- Problem Formulation}
  \begin{itemize}
    \item LSS $\Sigma$ consists of $N$ weakly coupled subsystems
      $\nth{\Sigma}$:
      \begin{align*} 
    \nth{\Sigma}:\quad \nth{\bfx}_{k+1} &= \nth{\bff} \left(
    \bfx_{k},\nth{\mu}_{k},\nth{\bfu}_{k} \right) +
    \nth{\bfF}(\nth{\mu}_k)\nth{\bfw}_k\\
    \nth{\bfy}_{k} &= \nth{\bfh} \left(
    \nth{\bfx}_{k},\nth{\mu}_{k}\right) + \nth{\bfH}(\nth{\mu}_k)\nth{\bfv}_k
      \end{align*}
 \item<2-> Assumptions:
\begin{enumerate}
  \item Noises $\nth{\bfw}_k$ and $\nth{\bfv}_k$ white and mutually independent
  \item Initial mode indices $\nth{\mu}_{0}$ and initial states $\nth{\bfx}_{0}$ are independent
  \item Mode indices  $\nth{\mu}_{k}$ are modelled as Markov processes\\ %$P(\nth{\mu}_{k+1}|\nth{\mu}_{k})$\\
      \alert{$\nth{\mu}_{k}=1$} -- fault-free behavior, \alert{$\nth{\mu}_{k}\in\{2,\ldots,\nth{M}\}$} -- faulty behavior
\end{enumerate}
  \end{itemize}
\end{frame}
\begin{frame}
  \frametitle{Dependent Faults}
  \begin{itemize}
    \item Model indices assumptions
      \begin{align*}
      \hspace*{-3em}
        \Pr(\bfmu_{k+1}|\bfmu_{k})=
        \begin{cases}
        \prod_{n=1}^N\Pr(\nth{\mu}_{k+1}|\nth{\mu}_{k}) & \text{\tiny in~Straka, Pun\v{c}och\'{a}\v{r} (2019)}\\
                                                          %&\text{or}\\
          \prod_{n=1}^N\Pr(\nth{\mu}_{k+1}|\bfmu_{k}) & \text{\tiny in~Straka, Pun\v{c}och\'{a}\v{r} (2020)}
        \end{cases}
      \end{align*}
    \item Insufficient in some cases as the faults in different subsystems tend to appear \alert{simultaneously} or even \alert{system-wide faults} may emerge.
    \item This behavior can be formally described as dependent faults.
    \item \alert{Goal of the paper: design AFD for dependent faults}
  \end{itemize}
\end{frame}
\begin{frame}
  \frametitle{Active Fault Diagnosis -- Problem Formulation}
  \begin{itemize}[<+->]
    \item AFD system:
      \begin{align*}
  \Delta:\quad
  \begin{bmatrix}
    \bfd_{k}\\
    \bfu_{k}
  \end{bmatrix} = 
  %\bfrho_{k}\left(\bfI_{0}^{k}\right) =
  \begin{bmatrix}
    \bfsigma_{k}\left( \bfI_0^k\right)\\
    \bfgamma_{k}\left(\bfI_{0}^{k}\right)
  \end{bmatrix},
\end{align*}
where $\bfI_0^k\triangleq[\bfy_{0}^{k},\bfu_{0}^{k-1}]$ is the information available at $k$
    \item AFD goal: generate $\bfd_k$ and $\bfu_k$ to minimize discounted
      criterion
      \begin{align*} 
        J(\bfsigma_{0}^{\infty},\bfgamma_{0}^{\infty}) = \lim_{F\to\infty} \mean \left\{ \sum_{k=0}^{F} \eta^{k} L^{\mathrm{d}}(\bfmu_{k},\bfd_{k})\right\}
      \end{align*}
      \item Example of cost function
\begin{align*}
L^{\mathrm{d}}(\bfmu_{k},\bfd_{k}) = \sum_{n=1}^{N} \nth{L}^{\mathrm{d}} \left(\nth{\mu}_{k},\nth{d}_{k} \right),
\end{align*}
where $\nth{L}^{\mathrm{d}}$ penalizes discrepancy between $\nth{\mu}_{k}$ and $\nth{d}_{k}$ (missed detections, false alerts, incorrect fault identifications)
  \end{itemize}
\end{frame}
\begin{frame}
  \frametitle{General Approach to AFD design}
  \begin{block}{AFD algorithm stages:}
    \begin{enumerate}
      \item off-line: design of the excitation input and decision generators
      \item on-line: state estimation and utilization of the generators
    \end{enumerate}
  \end{block}
  \uncover<2->{
   \textbf{Solution:} adopt dynamic programming and find Bellman function 
  }
  \uncover<3>{
  \begin{block}{Dim. of $\bfI_0^k$ increases with $k\rightarrow \infty$  $\implies$ problem reformulation:}
    \begin{itemize} 
      \item unknown system state $\bfs_k= [\bfx^{\T}_{k}, \mu_{k}]^{\T}$
      \item recast the \textbf{imperfect state
        information model} (unknown $\bfs_k$) to a \textbf{perfect state information model} with a new (known) \alert{information state $\bfxi_{k}$}
      \item $\bfxi_k$ consists of sufficient statistics related to $p(\bfs_k|\bfI_0^k)$ obtained by an estimator
    \end{itemize}
  \end{block}
}
\end{frame}
%\begin{frame}
  %\frametitle{Problem Reformulation}
  %\begin{center}
    %\begin{tikzpicture}
		%% Perfect state information model
      %\node [sys,fill=gray!40, minimum height = 3cm, minimum width = 4.8cm] (psim){};
      %\node [anchor=north west,inner sep=2pt,font=\footnotesize] at (psim.north west) {Perfect state information model};
		%% System
      %\node [wr,minimum height=1.1cm,minimum width=2.1cm] at ( $(psim.north)+(0,-0.9)$ ) (system) {System };
      %\node [anchor=south east,inner sep=1pt] at (system.south east) {$\bfs_{k} =[\bfx_{k},\mu_{k}]$};
		%% Estimator
      %\node [wr,minimum height=1.1cm,minimum width=2.1cm,below=0.3 of system, align=center] (estimator) {State \\ estimator};
		%% Input
      %\draw[very thick,->] ($(system.west)+(-2.0,0)$) node[above] {$\bfu_{k}$} -- (system.west);
      %\draw[very thick,->] ($(system.west)+(-0.75,0)$)  |- ($(estimator.west)+(0,-0.2)$);
		%% Output
      %\draw[very thick,->] ($(system.east)+(0,0)$)  -- ($(system.east)+(1,0)$) node[above, midway] {$\bfy_{k}$} -- ($(system.east)+(1,-0.7)$) -- ($(system.east)+(-2.5,-0.7)$) |- ($(estimator.west)+(0,0.2)$);
		%% Hyper state
      %\draw[very thick,->] (estimator.east)  -- ($(estimator.east)+(2.2,0)$) node[above] {$p(\bfs_{k}|\bfI_{0}^{k})$};
    %\end{tikzpicture}
  %\end{center}
  %State estimator -- Bayesian Recursive Relations
    %\begin{center}
      %\begin{tikzpicture}
	%\node[text width=2cm] at (0,0) (predictivek) {$p (\bfs_{k} \vert
    %\bfI_{0}^{k})$\\as $\bfxi_k$};
	%\node[text width=2cm] at (7,0) (predictivekp1) {$p (
		%\bfs_{k+1} \vert \bfI_{0}^{k+1} )$\\as $\bfxi_{k+1}$};
	%\draw[very thick,->] (predictivek) -- (predictivekp1) node[midway,above]
	%(m) {$\bfphi(\bfxi_k,\bfu_{k},\bfy_{k+1})$};
	%\draw[thick,->,color=red!60] (2,-3) -- (2,0) node[midway,rotate=90, above]
	%{$\textcolor{red!60}{p\left(\bfs_{k+1}|\bfs_{k},\bfu_{k}\right)}$};
	%\draw[thick,->,color=red!60] (2.7,-3) -- (2.7,0) node[midway,rotate=90, above]
	%{$\textcolor{red!60}{p\left(\bfy_{k+1}|\bfs_{k+1}\right)}$};
	%\draw[thick,->,color=blue!60] (4.3,-3) -- (4.3,0) node[midway,rotate=90, above]
	%{$\textcolor{blue!60}{\bfu_{k}}$};
	%\draw[thick,->,color=blue!60] (5,-3) -- (5,0) node[midway,rotate=90, above]
	%{$\textcolor{blue!60}{\bfy_{k+1}}$};
      %\end{tikzpicture}
    %\end{center}
%\end{frame}
\begin{frame}
  \frametitle{Design of AFD Systems -- Problem Reformulation -- Illustration}
  \vspace{-6.5cm}
  \hspace{1.7cm}
\begin{tikzpicture}[every node/.style={draw,thick,rounded
  corners},transform canvas={scale=0.9}]
\node[fill=blue!20,text width=4.5cm,align=center,draw=none] (tit_l) {\it Imperfect state information\\ problem};
\node[fill=blue!20,text width=4.5cm,align=center,draw=none,right=3.3cm of tit_l] (tit_r) {\it Perfect state information\\ problem};
    \node[below=6mm of tit_l,minimum width=3cm,label=95:{System model}] (sm_l)
    {\begin{minipage}{3cm}\vspace{-3ex}\begin{align*}
      \bfx_{k+1} &= f (\bfx_{k}, \mu_{k}, \bfu_{k} ) + \bfF_{\mu_{k}}\bfw_{k}, \\ \bfy_{k} &= \bfh (\bfx_{k}, \mu_{k})  + \bfH_{\mu_{k}}\bfv_{k},\\
      &P(\mu_{k+1}=j|\mu_{k}=i)
    \end{align*}\end{minipage}};

    \node[below=6mm of tit_r,minimum width=3cm,label=92:{System model and estimator}] (sm_r) {\begin{minipage}{3cm}\vspace{-2ex}\begin{align*} 
      \bfxi_{k+1} = \bfphi(\bfxi_{k}, \bfu_{k}, \bfy_{k+1}),
    \end{align*}\end{minipage}};

    \node[below=30mm of tit_l,xshift=1.2cm,minimum height=7mm] (e_l) {Estimate $\bfxi_k$};

    \node[below=40mm of tit_l,minimum width=3cm,minimum height=15mm,label=135:{Criterion}] (cr_l) {\begin{minipage}{3cm}\vspace{-2ex}\begin{align*}
      J = \lim_{F \to \infty} \mean \left\lbrace \sum^{F}_{k=0} \eta^{k}L^{d}(\mu_{k},d_{k}) \right\rbrace
    \end{align*}\end{minipage}};

    \node[below=40mm of tit_r,minimum width=3cm,minimum height=15mm,label=135:{Criterion}] (cr_r) {\begin{minipage}{3cm}\vspace{-3ex}\begin{align*}
      J = \lim_{F \to \infty} \mean \left\lbrace \sum^{F}_{k=0} \eta^{k}L^{d}(\bfxi_{k},d_{k}) \right\rbrace
    \end{align*}\end{minipage}};

    \node[below=5mm of cr_l,minimum width=3cm,label={[align=left]105:AFD}] (afd_l) {\begin{minipage}{3cm}\vspace{-2ex}\begin{align*}
      \left[ \begin{matrix}
  d_{k} \\ \bfu_{k}
      \end{matrix} \right]
      = \bfrho_{k}(\bfI^{k}_{0})
    \end{align*}\end{minipage}};
    \node[below=5mm of cr_r,minimum width=3cm,label={[align=left]105:AFD}] (afd_r) {\begin{minipage}{3cm}\vspace{-2ex}\begin{align*}
      \left[ \begin{matrix}
  d_{k} \\ \bfu_{k}
      \end{matrix} \right]
      = \bar{\bfrho}_{k}(\bfxi_{k})
    \end{align*}\end{minipage}};
    \draw[very thick,dashed,->] (sm_l) -- (cr_l);
    \draw[very thick,dashed,->] (cr_l) -- (afd_l);
    \draw[very thick,dashed,->] (e_l.north|-sm_l.south) -- (e_l);
    \draw[very thick,dashed,->] (sm_r) -- (cr_r);
    \draw[very thick,dashed,->] (cr_r) -- (afd_r);

\draw[fill=blue!40!green!25,draw=none] ($(tit_l.north east)+(0.3,0)$) rectangle ($(afd_r.south west)+(-1.2,0)$);
\node [rotate=90,right=1.1cm of e_l,text width=4cm,align=center,xshift=-1.3cm,draw=none] (pr) {\Large Problem\\ reformulation};
\node[below=3.8cm of pr,single arrow,rounded corners=0pt,minimum
width=1.5cm,minimum height=1.5cm,color=white,fill=white,xshift=-6.5mm] (a1) {\Large a};
\node[above=1.1cm of a1,single arrow,rounded corners=0pt,minimum width=1.5cm,minimum height=1.5cm,color=white,fill=white] (a2) {\Large a};
\node[above=6.1cm of a1,single arrow,rounded corners=0pt,minimum width=1.5cm,minimum height=1.5cm,color=white,fill=white] (a3) {\Large a};
    %\draw[fill=blue!40!green!25,draw=none] ($(tit_l.north east)+(0.5,0)$) rectangle ($(afd_r.south west)+(-1.8,0)$);
    %\node [rotate=90,right=1.9cm of e_l,text width=4cm,align=center,xshift=-1.3cm,draw=none] (pr) {\Large Problem\\ reformulation};
    %\node[below=3.8cm of pr,single arrow,rounded corners=0pt,minimum
    %width=1.5cm,minimum height=1.5cm,color=white,fill=white,xshift=-6.5mm] (a1) {\Large a};
    %\node[above=1.1cm of a1,single arrow,rounded corners=0pt,minimum width=1.5cm,minimum height=1.5cm,color=white,fill=white] (a2) {\Large a};
    %\node[above=5.9cm of a1,single arrow,rounded corners=0pt,minimum width=1.5cm,minimum height=1.5cm,color=white,fill=white] (a3) {\Large a};
  \end{tikzpicture}
\end{frame}
\begin{frame}
  \frametitle{Hierarchical AFD Architecture}
  \begin{itemize}
    \item off-line design will employ the computationally efficient decentralized architecture
    \item on-line estimation will work in a hierarchical architecture that has the ability to take into account the faults dependency and the coupling through~$\bfx_k$.
  \end{itemize}
\end{frame}
\begin{frame}
  \frametitle{Hierarchical AFD Architecture}
\begin{center}
    \begin{tikzpicture}[scale=0.7, every node/.style={scale=0.7}]
      \node (dex1) [wr] {$\ist{\Sigma}$};
      \node (dex2) [wr,right=1cm of dex1] {$\iind{\Sigma}$};
      \node (dex3) [wr,right=1cm of dex2] {$\iiird{\Sigma}$};

      \draw[latex-latex,thick] (dex1) -- (dex2);
      \draw[latex-latex,thick] (dex2) -- (dex3);
      \draw[latex-latex,thick] (dex1.north) to [out=10,in=170]  (dex3.north);
      \node (lde) at ($(dex1.north east)+(-1.8,-0.1)$) {LSS ($\Sigma$)};
      \node (defd1) [gr,below=7mm of dex1] {$\ist{\Delta}$};
      \node (defd2) [gr,below=7mm of dex2] {$\iind{\Delta}$};
      \node (defd3) [gr,below=7mm of dex3] {$\iiird{\Delta}$};
      \draw[whiVecArrow] (defd1) -- (dex1);
      \draw[whiVecArrow] (defd2) -- (dex2);
      \draw[whiVecArrow] (defd3) -- (dex3);
      \node (defd0) [pr,below=15mm of defd2] {$\zeroth{\Delta}$};
      \draw[grVecArrow] (defd1) -- (defd2);
      \draw[grVecArrow] (defd2) -- (defd3);
      \draw[grVecArrow1] (defd1.-45) to [out=-25,in=-155] (defd3.-135);
      \draw[purVecArrow1] (defd1) to [bend right=20] (defd0);
      \draw[purVecArrow] (defd2) -- (defd0);
      \draw[purVecArrow1] (defd3) to [bend left=20] (defd0);
      %\draw[-latex,thick] (defd3.-80) to [bend left=20]  (defd0.-10);
      %\draw[-latex,thick] (defd0.10) to [bend right=10]  (defd3.-100);
      %\draw[-latex,thick] (defd1.-80) to [bend right=10]  (defd0.170);
      %\draw[-latex,thick] (defd0.190) to [bend left=20]  (defd1.-100);
      %\draw[-latex,thick] (defd2.-80) to [bend left=5]  (defd0.80);
      %\draw[-latex,thick] (defd0.100) to [bend left=5]  (defd2.-100);
      \begin{pgfonlayer}{back}
        \draw[sys] ($(dex1.north west)+(-1.5,0.3)$) rectangle ($(dex3.south east)+(1.5,-.5)$);
      \end{pgfonlayer}
    \end{tikzpicture}
  \end{center}
  \vspace{-3ex}
  \begin{itemize}
    \item local nodes
      \begin{itemize}
        \item estimate the continuous state of the subsystems
        \item select the optimal excitations 
        \item send local likelihoods to the central node.
      \end{itemize}
    \item central node 
      \begin{itemize}
        \item generates the decisions
        \item submits the respective model probabilities to the local nodes
      \end{itemize}
  \end{itemize}
  %\begin{itemize}[<+->]
    %%\item Combination of centralized and decentralized AFDs
      %%\item State $\nth{\bar{\bfs}}_{k} \triangleq [\bfx^{\T}_{k}, \nth{\mu}_{k}]^{\T}$
      %\item To obtain its estimate, the nodes must communicate individual estimates (the measurement models are known only by the corresponding nodes).
      %\item The estimates are then fused (under unknown dependency).
      %\item To achieve feasible Bellman function calculation, the Bellman function from the decentralized architecture is used.
      %\item $\implies$ the input generator of decentralized architecture is used but it obtains better information by respecting the coupling
  %\end{itemize}
\end{frame}
%\begin{frame}
  %\frametitle{AFD for Distributed Architecture -- Algorithm }
%\begin{algorithm}[H]
%% input/ouput names
%\SetKwInOut{Input}{Input}
%\SetKwInOut{Output}{Output}

%\caption{Distributed AFD algorithm.}
%\Input{%
  %{$\nth{\bfy}_{k}$} -- observation of $\nth{\Sigma}$ \\
  %predictive estimate  of $\nth{\bfs}_{k}$ given by
  %$p(\nth{\bfs}_k\vert\nth{\bfI}_0^{k-1})$
%}
%\Output{%
  %{${\nth{d}_k}$} -- decision about the current model $\nth{\mu}_k$\\
  %{${\nth{\bfu}_k}$} -- excitation input for  the current model
  %$\nth{\Sigma}_k$\\
  %predictive estimate  of $\nth{\bfs}_{k+1}$ given by
  %$p(\nth{\bfs}_{k+1}\vert\nth{\bfI}_0^{k})$
%}

  %\BlankLine % blank line for spacing
  
  %% start of the pseudocode
  %\textbf{Filtering:} Update the estimate $\nth{\bfs}_{k}$
  %using $\nth{\bfy}_{k}$  to obtain
  %$p(\nth{\bfs}_{k}\vert\bfI_0^{k-1},\nth{\bfy}_{k},\nth{\bfu}_{k-1})$.

  %\textbf{Merging:} Merge the filtering PDF obtained in the previous step
  %using GPB1 to $\nth{M}$ terms related to the models
  %$\nth{\mu}_{k}$ and express the merged PDF using
  %$\nth{\bfxi}_{k}$.

  %\textbf{Communication:} Obtain the information states $\ith{\bfxi}_k$, $i\neq n$ from other nodes.

  %\textbf{Fusion:} Obtain the mixture $p(\bfs_k\vert\bfI_0^k)$ and statistics $\bfxi_k$ related to $\Sigma$.

  %\textbf{Decision generation:} Generate the
  %decision $\nth{d}_k$ using $\bfxi_k$.

  %\textbf{Input generation:} Generate the
  %excitation input $\nth{\bfu}_k$ from $\nth{\bfxi}_k$.

  %\textbf{Prediction:} Calculate $\nth{\bfs}_{k+1}$ prediction estimate using $\nth{\bfu}_k$ and $\bfxi_k$.
%\end{algorithm}
%\end{frame}
\begin{frame}
  \frametitle{Numerical Illustration -- System Specification}
  Two weakly coupled multiple-model linear subsystems $\ith{\Sigma},\ i=1,2$
\begin{align*}
  \ith{\Sigma}:\quad \ith{\bfx}_{k+1} &= \ith{\bfA}_{(\ith{\mu}_{k})} \bfx_{k} +
  \ith{B}_{(\ith{\mu}_{k})}\ith{u}_{k} +
  \ith{G}_{(\ith{\mu}_{k})}\ith{w}_{k},\\
  \ith{z}_k&=\ith{C}_{(\ith{\mu}_{k})}\ith{\bfx}_{k}+\ith{H}_{(\ith{\mu}_{k})}\ith{v}_k\\
  %\iind{\Sigma}: \iind{\bfx}_{k+1} &=  \iind{\bfA}_{(\iind{\mu}_{k})} \bfx_{k} +
  %\iind{B}_{(\iind{\mu}_{k})}\iind{u}_{k} +
  %\ist{G}_{(\iind{\mu}_{k})}\iind{w}_{k},\\
  %\iind{z}_k&=\iind{C}_{(\iind{\mu}_{k})}\iind{\bfx}_{k}+\iind{H}_{(\iind{\mu}_{k})}\iind{v}_k
\end{align*}
where both subsystems have two modes with the following matrices
    \begin{center}
      \scalebox{0.7}{\parbox{\textwidth}{%
\begin{align*}
  \ist{\Sigma}:\quad
  \begin{cases}
  \ist{\bfA}_{(1)} = [0.98\ 0.01], \ist{B}_{(1)} = 0.01,\ist{G}_{(1)}=\sqrt{0.003}, \ist{C}_{(1)}=1, \ist{H}_{(1)}=0.01,\\
  \ist{\bfA}_{(2)}=[0.92\ 0.03], \ist{B}_{(2)}=0.08,\ist{G}_{(2)}=\sqrt{0.003}, \ist{C}_{(2)}=1, \ist{H}_{(2)}=0.01,\\
  \end{cases}\\
  \iind{\Sigma}:\quad
  \begin{cases}
  \iind{\bfA}_{(1)} = [0.02\ 0.93],\iind{B}_{(1)}= 0.07,\iind{G}_{(1)}=\sqrt{0.002}, \iind{C}_{(1)}=1, \iind{H}_{(1)}=0.01,\\
  \iind{\bfA}_{(2)} = [0.01\ 0.93], \iind{B}_{(2)} = 0.09,\iind{G}_{(2)}=\sqrt{0.002} \iind{C}_{(2)}=1, \iind{H}_{(2)}=0.01.
\end{cases}
\end{align*}}}
\end{center}
\end{frame}
\begin{frame}
  \frametitle{Numerical Illustration --  Transition probabilities}
\begin{center}
  \begin{tabular}{rlllll}\toprule
     & \multicolumn{4}{c}{$(\ist{\mu}_{k},\iind{\mu}_{k})$}\\
    \cmidrule{2-5}
    $(\ist{\mu}_{k+1},\iind{\mu}_{k+1})$ & (1,1) &  (1,2) & (2,1) & (2,2)\\\midrule 
    (1,1) & 0.95 & 0.04 & 0.04 & 0.01 \\
    (1,2) & 0.02 & 0.80 & 0.01 & 0.02 \\
    (2,1) & 0.02 & 0.01 & 0.80 & 0.02 \\
    (2,2) & 0.01 & 0.15 & 0.15 & 0.95 \\\bottomrule
  \end{tabular}
\end{center}
%\begin{table}
  %\definecolor{Gray}{gray}{0.85}
  %\newcolumntype{a}{>{\columncolor{Gray}}c}
  %\definecolor{Gray1}{gray}{0.85}
  %\newcolumntype{b}{>{\columncolor{Gray1}}c}
  %\extrarowheight=\aboverulesep
  %\addtolength{\extrarowheight}{\belowrulesep}
%\aboverulesep=0pt
%\belowrulesep=0pt
 %\begin{adjustbox}{width=.5\textwidth, totalheight=\textheight-2\baselineskip,keepaspectratio}
  %\begin{tabular}{llllll}\toprule
    %%\multicolumn{5}{c}{$\ist{\Sigma}$} & \phantom{a} & \multicolumn{5}{c}{$\iind{\Sigma}$} \\\cmidrule{1-5}\cmidrule{7-11}
    %&&\multicolumn{4}{c}{$[\ist{\mu}_{k},\iind{\mu}_k]$}\\
     %\cmidrule{3-6}
    %&& [1,1] &[1,2] & [2,1] & [2,2]\\
    %\midrule
    %\multirow{2}{*}{$\ist{\mu}_{k+1}$}&
    %1 & 0.99 & 0.9 & 0.1 & 0.02\\
      %&2 & 0.01 & 0.1 & 0.9 & 0.98\\
    %%\midrule
      %%&\multicolumn{4}{c}{$[\ist{\mu}_{k},\iind{\mu}_k]$}\\
      %%\cmidrule{2-5}
      %\multirow{2}{*}{$\iind{\mu}_{k+1}$}%& [1,1] &[1,2] &[2,1] & [2,2]\\ 
       %&1 & 0.99 &  0.05 & 0.9 & 0.01\\
       %&2 & 0.01 &  0.95 & 0.1 & 0.99\\
    %\bottomrule
  %\end{tabular}
%\end{adjustbox}
%\begin{itemize}
  %\item $p_{\ist{\bfw_{k}}}(\ist{\bfw_{k}})=p_{\iind{\bfw_{k}}}(\iind{\bfw_{k}})=\mcN\{0,1\}$
  %\item $p_{\bfx_{0}}(\bfx_0)= \mcN\{\bfZero,0.01\cdot\bfIdent\}$,
    %$P(\bfmu_{0}=[1\ 1]^{\T})=1$
  %\item Admissible inputs of subsystems:  $\ist{\mcU}=\iind{\mcU} = \{-1,0,1\}$ 
  %\item Detection cost function: $L^{\mathrm{d}}(\nth{\mu}_{k},\nth{d}_{k}) =
    %1-\delta_{\nth{\mu}_{k},\nth{d}_{k}}$ ($\eta=0.98$).
  %\item Value iteration over grid of $732\, 050$ discrete states in
    %%$[-1:0.2:1]\times[-1:0.2:1]\times[0:2\cdot10^{-5}:2\cdot10^{-4}]\times[0:2\cdot10^{-5}:2\cdot10^{-4}]\times[0:0.02:1]$
    %$[-1,1]\times[-1,1]\times[0,2\cdot10^{-4}]\times[0,2\cdot10^{-4}]\times[0,1]$
%\end{itemize}
%\end{table}
\end{frame}
%\begin{frame}
  %\frametitle{Numerical Illustration  -- Architectures}
  %\begin{itemize}
    %\item decentralized architecture: ignoring  the coupling, two independent
      %information state -- 5 dimensions
    %\item distributed architecture: input signal generator of decentralized
      %architecture, information state -- 11 dimensions
      %\item centralized architecture: information state -- 22 dimensions
        %$\implies$ infeasible Bellman function calculation
  %\end{itemize}
%\end{frame}
\begin{frame}
  \frametitle{Numerical illustration -- Performance Evaluation}
  \begin{itemize}
    \item  $10^5$ Monte Carlo simulations
    \item Finite time horizon: $F=400$
    \item PFD -- passive fault diagnosis with $u_k=\sign(\sin(k))$ 
    \item Results:
  \end{itemize}
  \begin{center}
    \begin{tabular}{llS[table-format=2.2]S[table-format=2.2]l}\toprule
      &$\hat{J}$& $P_{\text{MD}}$ & $P_{\text{FA}}$ & $T_{\text{on-line}}$\\\midrule
      decentralized AFD     & 2.876 & 12.94\si{\percent} & 10.08\si{\percent} &0.150 \si{\second}\\
      distributed AFD     & 1.852 & 2.06\si{\percent} & 2.56\si{\percent} &0.421 \si{\second}\\
      hierarchical AFD      & 1.691 & 1.97\si{\percent} & 2.30\si{\percent} &0.487 \si{\second}\\\bottomrule
    \end{tabular} 
  \end{center}
\end{frame}
\begin{frame}
  \frametitle{Conclusion}
  \begin{itemize}
    \item The paper focused on the active fault diagnosis of a large scale stochastic system with coupled faults.
    \item The system was represented using multiple models, which describe the fault-free and faulty behavior.
    \item The off-line design was carried out in a decentralized architecture to achieve computational and storage tractability.
    \item The on-line estimation used a hierarchical architecture.
    \item The numerical example demonstrated improved detection quality of the proposed hierarchical architecture.
  \end{itemize}
\end{frame}
\frame{\titlepage}
\end{document}

